% --- Archon physics formalization source begin ---
% archon:physics
% archon:covers PhyXMiniProblems/problem_phyx_mini_0853.lean
% archon:source-report reports/phyx_mini/problem_phyx_mini_0853.source.json

\chapter{Physics problem phyx\_mini\_0853}
\label{ch:PhyXMiniProblems_problem_phyx_mini_0853}

\paragraph{Problem source.}
\#\# Physical scenario

Three electrons are spaced 1.0 mm apart along a vertical line. The outer two electrons are fixed in position.

\#\# Question

If the center electron is displaced horizontally by a small distance, what will its speed be when it is very far away?

\#\# Answer choices

- A: A: 1500m/s

- B: B: 1400m/s

- C: C: 1250m/s

- D: D: 1000m/s

\#\# Figure caption (auxiliary; use the image as primary evidence)

The image is a diagram illustrating an interaction between three charged particles. Here's a breakdown of its components and relationships:

1. **Before Section:**
   - Three negatively charged particles are shown vertically aligned and labeled as 1, 2, and 3.
   - Distances between charges are indicated:
     - \textbackslash{}((r\_\{12\})\_i = 1.0 \textbackslash{}, \textbackslash{}text\{mm\}\textbackslash{}) between particle 1 and 2.
     - \textbackslash{}((r\_\{23\})\_i = 1.0 \textbackslash{}, \textbackslash{}text\{mm\}\textbackslash{}) between particle 2 and 3.
   - There's a velocity labeled \textbackslash{}(v\_i = 0\textbackslash{}) next to particle 2, indicating initial velocity is zero.

2. **After Section:**
   - Particle 2, now has a velocity labeled \textbackslash{}(v\_f\textbackslash{}), represented by a green arrow pointing to the right.
   - The distances \textbackslash{}((r\_\{12\})\_f\textbackslash{}) and \textbackslash{}((r\_\{23\})\_f\textbackslash{}) are shown with dashed lines extending to infinity, implying particles 1 and 3 have moved far apart from particle 2.

3. **Overall Layout:**
   - The image is divided into "Before" and "After" sections, showing the initial and final states of the particles.
   - The charges are consistently negative, and their interactions likely relate to electrostatic forces.
   - Textual details include descriptions of distances and velocities, supporting the conceptual visualization of particle movement.

\#\# Dataset metadata

Electrostatics; Physical Model Grounding Reasoning, Multi-Formula Reasoning

\paragraph{Recorded answer/context.}
D: D: 1000m/s

\paragraph{Figure/image path.}
/root/proposal\_for\_physic/science-mango/phyx\_mini\_run/phyx\_data/test\_image/853.png

\paragraph{Formalization target.}
create a compiling Lean file with sorry bodies at `PhyXMiniProblems/problem\_phyx\_mini\_0853.lean`. The Lean declarations must preserve the physical quantities, dimensions or dimensional roles, figure labels, governing-law hypotheses, and final relation expressed by this problem.
Use Mathlib/Physlib names found through LeanExplore where available. If a physics API is missing, introduce faithful local abstractions rather than scalar placeholder aliases.

\begin{theorem}[Physics formalization target]
\label{thm:physics:phyx_mini_0853:target}
\lean{PhyXMiniProblems.ProblemPhyXMini0853.problem_phyx_mini_0853}
\uses{def:physics:phyx-mini-0853:phyxminiproblems-problemphyxmini0853-lengthinmeters, def:physics:phyx-mini-0853:phyxminiproblems-problemphyxmini0853-massinkilograms, def:physics:phyx-mini-0853:phyxminiproblems-problemphyxmini0853-chargeincoulombs, def:physics:phyx-mini-0853:phyxminiproblems-problemphyxmini0853-speedinmeterspersecond, def:physics:phyx-mini-0853:phyxminiproblems-problemphyxmini0853-coulombconstantinjoulemeterspercoulombsquared, def:physics:phyx-mini-0853:phyxminiproblems-problemphyxmini0853-centerouterpair, def:physics:phyx-mini-0853:phyxminiproblems-problemphyxmini0853-threeelectronescapesetup, def:physics:phyx-mini-0853:phyxminiproblems-problemphyxmini0853-matchesthreeelectronscenario, def:physics:phyx-mini-0853:phyxminiproblems-problemphyxmini0853-matchesprimarythreeelectronfigure, def:physics:phyx-mini-0853:phyxminiproblems-problemphyxmini0853-hasdepictedandreleasegeometry, def:physics:phyx-mini-0853:phyxminiproblems-problemphyxmini0853-hassmallphysicalperturbation, def:physics:phyx-mini-0853:phyxminiproblems-problemphyxmini0853-useselectronandvacuumreferencedata, def:physics:phyx-mini-0853:phyxminiproblems-problemphyxmini0853-satisfiescoulombescapeandenergylaws, def:physics:phyx-mini-0853:phyxminiproblems-problemphyxmini0853-isuniqueclosestdisplayedspeed}
The assigned autoformalize agent should translate this physics problem into Lean declarations in the covered file, with theorem and lemma proof bodies written as `by sorry`.
\end{theorem}
\begin{proof}
This is an autoformalization task, not a proof task. Produce statements that can later be proved by the physics prover without weakening the physical model.
\end{proof}

% --- Archon declaration topology begin ---
\section{Declaration topology}
\begin{definition}[coulomb Constant Dimension]
  \label{def:physics:phyx-mini-0853:phyxminiproblems-problemphyxmini0853-coulombconstantdimension}
  \lean{PhyXMiniProblems.ProblemPhyXMini0853.coulombConstantDimension}
  The physical dimension of Coulomb's constant, equivalently J m / C²
\end{definition}
\begin{proof}
  This is the declared model definition of coulomb Constant Dimension.
\end{proof}
\begin{definition}[Length Quantity]
  \label{def:physics:phyx-mini-0853:phyxminiproblems-problemphyxmini0853-lengthquantity}
  \lean{PhyXMiniProblems.ProblemPhyXMini0853.LengthQuantity}
  A nonnegative, unit-independent physical length
\end{definition}
\begin{proof}
  This is the declared model definition of Length Quantity.
\end{proof}
\begin{definition}[Planar Position Quantity]
  \label{def:physics:phyx-mini-0853:phyxminiproblems-problemphyxmini0853-planarpositionquantity}
  \lean{PhyXMiniProblems.ProblemPhyXMini0853.PlanarPositionQuantity}
  A unit-independent signed planar position
\end{definition}
\begin{proof}
  This is the declared model definition of Planar Position Quantity.
\end{proof}
\begin{definition}[Mass Quantity]
  \label{def:physics:phyx-mini-0853:phyxminiproblems-problemphyxmini0853-massquantity}
  \lean{PhyXMiniProblems.ProblemPhyXMini0853.MassQuantity}
  A positive physical mass, represented independently of any unit choice
\end{definition}
\begin{proof}
  This is the declared model definition of Mass Quantity.
\end{proof}
\begin{definition}[Signed Charge Quantity]
  \label{def:physics:phyx-mini-0853:phyxminiproblems-problemphyxmini0853-signedchargequantity}
  \lean{PhyXMiniProblems.ProblemPhyXMini0853.SignedChargeQuantity}
  A signed, unit-independent electric charge
\end{definition}
\begin{proof}
  This is the declared model definition of Signed Charge Quantity.
\end{proof}
\begin{definition}[Speed Quantity]
  \label{def:physics:phyx-mini-0853:phyxminiproblems-problemphyxmini0853-speedquantity}
  \lean{PhyXMiniProblems.ProblemPhyXMini0853.SpeedQuantity}
  A nonnegative physical speed
\end{definition}
\begin{proof}
  This is the declared model definition of Speed Quantity.
\end{proof}
\begin{definition}[Energy Quantity]
  \label{def:physics:phyx-mini-0853:phyxminiproblems-problemphyxmini0853-energyquantity}
  \lean{PhyXMiniProblems.ProblemPhyXMini0853.EnergyQuantity}
  Signed physical kinetic or electrostatic potential energy
\end{definition}
\begin{proof}
  This is the declared model definition of Energy Quantity.
\end{proof}
\begin{definition}[Coulomb Constant Quantity]
  \label{def:physics:phyx-mini-0853:phyxminiproblems-problemphyxmini0853-coulombconstantquantity}
  \lean{PhyXMiniProblems.ProblemPhyXMini0853.CoulombConstantQuantity}
  \uses{def:physics:phyx-mini-0853:phyxminiproblems-problemphyxmini0853-coulombconstantdimension}
  A nonnegative, dimensionful value of Coulomb's constant
\end{definition}
\begin{proof}
  This is the declared model definition of Coulomb Constant Quantity.
\end{proof}
\begin{definition}[x Axis]
  \label{def:physics:phyx-mini-0853:phyxminiproblems-problemphyxmini0853-xaxis}
  \lean{PhyXMiniProblems.ProblemPhyXMini0853.xAxis}
  Coordinate 0, horizontal and positive toward the right of the image
\end{definition}
\begin{proof}
  This is the declared model definition of x Axis.
\end{proof}
\begin{definition}[y Axis]
  \label{def:physics:phyx-mini-0853:phyxminiproblems-problemphyxmini0853-yaxis}
  \lean{PhyXMiniProblems.ProblemPhyXMini0853.yAxis}
  Coordinate 1, vertical and positive toward the top of the image
\end{definition}
\begin{proof}
  This is the declared model definition of y Axis.
\end{proof}
\begin{definition}[length Readout]
  \label{def:physics:phyx-mini-0853:phyxminiproblems-problemphyxmini0853-lengthreadout}
  \lean{PhyXMiniProblems.ProblemPhyXMini0853.lengthReadout}
  \uses{def:physics:phyx-mini-0853:phyxminiproblems-problemphyxmini0853-lengthquantity}
  Read a physical length in a selected named unit
\end{definition}
\begin{proof}
  This is the declared model definition of length Readout.
\end{proof}
\begin{definition}[length In Meters]
  \label{def:physics:phyx-mini-0853:phyxminiproblems-problemphyxmini0853-lengthinmeters}
  \lean{PhyXMiniProblems.ProblemPhyXMini0853.lengthInMeters}
  \uses{def:physics:phyx-mini-0853:phyxminiproblems-problemphyxmini0853-lengthquantity, def:physics:phyx-mini-0853:phyxminiproblems-problemphyxmini0853-lengthreadout}
  Read a physical length in coherent-SI metres
\end{definition}
\begin{proof}
  This is the declared model definition of length In Meters.
\end{proof}
\begin{definition}[length In Millimeters]
  \label{def:physics:phyx-mini-0853:phyxminiproblems-problemphyxmini0853-lengthinmillimeters}
  \lean{PhyXMiniProblems.ProblemPhyXMini0853.lengthInMillimeters}
  \uses{def:physics:phyx-mini-0853:phyxminiproblems-problemphyxmini0853-lengthquantity, def:physics:phyx-mini-0853:phyxminiproblems-problemphyxmini0853-lengthinmeters}
  Read a physical length in millimetres, as used by the two figure labels
\end{definition}
\begin{proof}
  This is the declared model definition of length In Millimeters.
\end{proof}
\begin{definition}[position In Meters]
  \label{def:physics:phyx-mini-0853:phyxminiproblems-problemphyxmini0853-positioninmeters}
  \lean{PhyXMiniProblems.ProblemPhyXMini0853.positionInMeters}
  \uses{def:physics:phyx-mini-0853:phyxminiproblems-problemphyxmini0853-planarpositionquantity}
  Read a planar physical position as Cartesian coordinates in metres
\end{definition}
\begin{proof}
  This is the declared model definition of position In Meters.
\end{proof}
\begin{definition}[mass In Kilograms]
  \label{def:physics:phyx-mini-0853:phyxminiproblems-problemphyxmini0853-massinkilograms}
  \lean{PhyXMiniProblems.ProblemPhyXMini0853.massInKilograms}
  \uses{def:physics:phyx-mini-0853:phyxminiproblems-problemphyxmini0853-massquantity}
  Read a physical mass in coherent-SI kilograms
\end{definition}
\begin{proof}
  This is the declared model definition of mass In Kilograms.
\end{proof}
\begin{definition}[charge In Coulombs]
  \label{def:physics:phyx-mini-0853:phyxminiproblems-problemphyxmini0853-chargeincoulombs}
  \lean{PhyXMiniProblems.ProblemPhyXMini0853.chargeInCoulombs}
  \uses{def:physics:phyx-mini-0853:phyxminiproblems-problemphyxmini0853-signedchargequantity}
  Read a signed physical charge in coherent-SI coulombs
\end{definition}
\begin{proof}
  This is the declared model definition of charge In Coulombs.
\end{proof}
\begin{definition}[speed In Meters Per Second]
  \label{def:physics:phyx-mini-0853:phyxminiproblems-problemphyxmini0853-speedinmeterspersecond}
  \lean{PhyXMiniProblems.ProblemPhyXMini0853.speedInMetersPerSecond}
  \uses{def:physics:phyx-mini-0853:phyxminiproblems-problemphyxmini0853-speedquantity}
  Read a speed in coherent-SI metres per second
\end{definition}
\begin{proof}
  This is the declared model definition of speed In Meters Per Second.
\end{proof}
\begin{definition}[energy In Joules]
  \label{def:physics:phyx-mini-0853:phyxminiproblems-problemphyxmini0853-energyinjoules}
  \lean{PhyXMiniProblems.ProblemPhyXMini0853.energyInJoules}
  \uses{def:physics:phyx-mini-0853:phyxminiproblems-problemphyxmini0853-energyquantity}
  Read a physical energy in coherent-SI joules
\end{definition}
\begin{proof}
  This is the declared model definition of energy In Joules.
\end{proof}
\begin{definition}[coulomb Constant In Joule Meters Per Coulomb Squared]
  \label{def:physics:phyx-mini-0853:phyxminiproblems-problemphyxmini0853-coulombconstantinjoulemeterspercoulombsquared}
  \lean{PhyXMiniProblems.ProblemPhyXMini0853.coulombConstantInJouleMetersPerCoulombSquared}
  \uses{def:physics:phyx-mini-0853:phyxminiproblems-problemphyxmini0853-coulombconstantquantity}
  Read Coulomb's constant in coherent SI, J m / C²
\end{definition}
\begin{proof}
  This is the declared model definition of coulomb Constant In Joule Meters Per Coulomb Squared.
\end{proof}
\begin{definition}[elementary Charge In Coulombs]
  \label{def:physics:phyx-mini-0853:phyxminiproblems-problemphyxmini0853-elementarychargeincoulombs}
  \lean{PhyXMiniProblems.ProblemPhyXMini0853.elementaryChargeInCoulombs}
  Physlib's elementary-charge unit expressed as a number of coulombs
\end{definition}
\begin{proof}
  This is the declared model definition of elementary Charge In Coulombs.
\end{proof}
\begin{definition}[Electron Label]
  \label{def:physics:phyx-mini-0853:phyxminiproblems-problemphyxmini0853-electronlabel}
  \lean{PhyXMiniProblems.ProblemPhyXMini0853.ElectronLabel}
  The particle numbers printed in the supplied diagram
\end{definition}
\begin{proof}
  This is the declared model definition of Electron Label.
\end{proof}
\begin{definition}[Outer Electron]
  \label{def:physics:phyx-mini-0853:phyxminiproblems-problemphyxmini0853-outerelectron}
  \lean{PhyXMiniProblems.ProblemPhyXMini0853.OuterElectron}
  The two fixed outer electrons
\end{definition}
\begin{proof}
  This is the declared model definition of Outer Electron.
\end{proof}
\begin{definition}[electron Label]
  \label{def:physics:phyx-mini-0853:phyxminiproblems-problemphyxmini0853-outerelectron-electronlabel}
  \lean{PhyXMiniProblems.ProblemPhyXMini0853.OuterElectron.electronLabel}
  \uses{def:physics:phyx-mini-0853:phyxminiproblems-problemphyxmini0853-electronlabel, def:physics:phyx-mini-0853:phyxminiproblems-problemphyxmini0853-outerelectron}
  The particle number corresponding to an outer-electron label
\end{definition}
\begin{proof}
  This is the declared model definition of electron Label.
\end{proof}
\begin{definition}[Center Outer Pair]
  \label{def:physics:phyx-mini-0853:phyxminiproblems-problemphyxmini0853-centerouterpair}
  \lean{PhyXMiniProblems.ProblemPhyXMini0853.CenterOuterPair}
  The two pair-distance labels displayed in the figure
\end{definition}
\begin{proof}
  This is the declared model definition of Center Outer Pair.
\end{proof}
\begin{definition}[outer Electron]
  \label{def:physics:phyx-mini-0853:phyxminiproblems-problemphyxmini0853-centerouterpair-outerelectron}
  \lean{PhyXMiniProblems.ProblemPhyXMini0853.CenterOuterPair.outerElectron}
  \uses{def:physics:phyx-mini-0853:phyxminiproblems-problemphyxmini0853-outerelectron, def:physics:phyx-mini-0853:phyxminiproblems-problemphyxmini0853-centerouterpair}
  The fixed outer electron occurring in a displayed pair
\end{definition}
\begin{proof}
  This is the declared model definition of outer Electron.
\end{proof}
\begin{definition}[Escape State]
  \label{def:physics:phyx-mini-0853:phyxminiproblems-problemphyxmini0853-escapestate}
  \lean{PhyXMiniProblems.ProblemPhyXMini0853.EscapeState}
  The two energy-balance states: release and asymptotically far away
\end{definition}
\begin{proof}
  This is the declared model definition of Escape State.
\end{proof}
\begin{definition}[Diagram Panel]
  \label{def:physics:phyx-mini-0853:phyxminiproblems-problemphyxmini0853-diagrampanel}
  \lean{PhyXMiniProblems.ProblemPhyXMini0853.DiagramPanel}
  The headings written over the two halves of the supplied diagram
\end{definition}
\begin{proof}
  This is the declared model definition of Diagram Panel.
\end{proof}
\begin{definition}[Horizontal Direction]
  \label{def:physics:phyx-mini-0853:phyxminiproblems-problemphyxmini0853-horizontaldirection}
  \lean{PhyXMiniProblems.ProblemPhyXMini0853.HorizontalDirection}
  The only horizontal direction used by the perturbation and velocity arrow
\end{definition}
\begin{proof}
  This is the declared model definition of Horizontal Direction.
\end{proof}
\begin{definition}[Separation Regime]
  \label{def:physics:phyx-mini-0853:phyxminiproblems-problemphyxmini0853-separationregime}
  \lean{PhyXMiniProblems.ProblemPhyXMini0853.SeparationRegime}
  Whether a depicted pair distance is finite or asymptotically infinite
\end{definition}
\begin{proof}
  This is the declared model definition of Separation Regime.
\end{proof}
\begin{definition}[Particle Species]
  \label{def:physics:phyx-mini-0853:phyxminiproblems-problemphyxmini0853-particlespecies}
  \lean{PhyXMiniProblems.ProblemPhyXMini0853.ParticleSpecies}
  Particle species distinguished by the problem statement
\end{definition}
\begin{proof}
  This is the declared model definition of Particle Species.
\end{proof}
\begin{definition}[Charge Sign]
  \label{def:physics:phyx-mini-0853:phyxminiproblems-problemphyxmini0853-chargesign}
  \lean{PhyXMiniProblems.ProblemPhyXMini0853.ChargeSign}
  Electric-charge signs visible inside the three circles
\end{definition}
\begin{proof}
  This is the declared model definition of Charge Sign.
\end{proof}
\begin{definition}[Three Electron Escape Figure]
  \label{def:physics:phyx-mini-0853:phyxminiproblems-problemphyxmini0853-threeelectronescapefigure}
  \lean{PhyXMiniProblems.ProblemPhyXMini0853.ThreeElectronEscapeFigure}
  \uses{def:physics:phyx-mini-0853:phyxminiproblems-problemphyxmini0853-lengthquantity, def:physics:phyx-mini-0853:phyxminiproblems-problemphyxmini0853-electronlabel, def:physics:phyx-mini-0853:phyxminiproblems-problemphyxmini0853-centerouterpair, def:physics:phyx-mini-0853:phyxminiproblems-problemphyxmini0853-diagrampanel, def:physics:phyx-mini-0853:phyxminiproblems-problemphyxmini0853-horizontaldirection, def:physics:phyx-mini-0853:phyxminiproblems-problemphyxmini0853-separationregime, def:physics:phyx-mini-0853:phyxminiproblems-problemphyxmini0853-chargesign}
  Literal content of image 853.png. The initial distance markers are physical lengths; numerical millimetre readouts are supplied separately
\end{definition}
\begin{proof}
  This is the declared model definition of Three Electron Escape Figure.
\end{proof}
\begin{definition}[Three Electron Escape Setup]
  \label{def:physics:phyx-mini-0853:phyxminiproblems-problemphyxmini0853-threeelectronescapesetup}
  \lean{PhyXMiniProblems.ProblemPhyXMini0853.ThreeElectronEscapeSetup}
  \uses{def:physics:phyx-mini-0853:phyxminiproblems-problemphyxmini0853-lengthquantity, def:physics:phyx-mini-0853:phyxminiproblems-problemphyxmini0853-planarpositionquantity, def:physics:phyx-mini-0853:phyxminiproblems-problemphyxmini0853-massquantity, def:physics:phyx-mini-0853:phyxminiproblems-problemphyxmini0853-signedchargequantity, def:physics:phyx-mini-0853:phyxminiproblems-problemphyxmini0853-speedquantity, def:physics:phyx-mini-0853:phyxminiproblems-problemphyxmini0853-energyquantity, def:physics:phyx-mini-0853:phyxminiproblems-problemphyxmini0853-coulombconstantquantity, def:physics:phyx-mini-0853:phyxminiproblems-problemphyxmini0853-electronlabel, def:physics:phyx-mini-0853:phyxminiproblems-problemphyxmini0853-outerelectron, def:physics:phyx-mini-0853:phyxminiproblems-problemphyxmini0853-centerouterpair, def:physics:phyx-mini-0853:phyxminiproblems-problemphyxmini0853-escapestate, def:physics:phyx-mini-0853:phyxminiproblems-problemphyxmini0853-horizontaldirection, def:physics:phyx-mini-0853:phyxminiproblems-problemphyxmini0853-particlespecies, def:physics:phyx-mini-0853:phyxminiproblems-problemphyxmini0853-threeelectronescapefigure}
  The unknown final speed and all intermediate energies are independent fields. No field is defined using the recorded answer 1000 m/s
\end{definition}
\begin{proof}
  This is the declared model definition of Three Electron Escape Setup.
\end{proof}
\begin{definition}[Matches Three Electron Scenario]
  \label{def:physics:phyx-mini-0853:phyxminiproblems-problemphyxmini0853-matchesthreeelectronscenario}
  \lean{PhyXMiniProblems.ProblemPhyXMini0853.MatchesThreeElectronScenario}
  \uses{def:physics:phyx-mini-0853:phyxminiproblems-problemphyxmini0853-speedinmeterspersecond, def:physics:phyx-mini-0853:phyxminiproblems-problemphyxmini0853-threeelectronescapesetup}
  Qualitative content stated in the problem, excluding the requested speed
\end{definition}
\begin{proof}
  This is the declared model definition of Matches Three Electron Scenario.
\end{proof}
\begin{definition}[Matches Primary Three Electron Figure]
  \label{def:physics:phyx-mini-0853:phyxminiproblems-problemphyxmini0853-matchesprimarythreeelectronfigure}
  \lean{PhyXMiniProblems.ProblemPhyXMini0853.MatchesPrimaryThreeElectronFigure}
  \uses{def:physics:phyx-mini-0853:phyxminiproblems-problemphyxmini0853-lengthinmillimeters, def:physics:phyx-mini-0853:phyxminiproblems-problemphyxmini0853-threeelectronescapesetup}
  All unambiguous words, labels, signs, and arrows visible in image 853.png
\end{definition}
\begin{proof}
  This is the declared model definition of Matches Primary Three Electron Figure.
\end{proof}
\begin{definition}[Has Depicted And Release Geometry]
  \label{def:physics:phyx-mini-0853:phyxminiproblems-problemphyxmini0853-hasdepictedandreleasegeometry}
  \lean{PhyXMiniProblems.ProblemPhyXMini0853.HasDepictedAndReleaseGeometry}
  \uses{def:physics:phyx-mini-0853:phyxminiproblems-problemphyxmini0853-xaxis, def:physics:phyx-mini-0853:phyxminiproblems-problemphyxmini0853-yaxis, def:physics:phyx-mini-0853:phyxminiproblems-problemphyxmini0853-lengthinmeters, def:physics:phyx-mini-0853:phyxminiproblems-problemphyxmini0853-positioninmeters, def:physics:phyx-mini-0853:phyxminiproblems-problemphyxmini0853-outerelectron, def:physics:phyx-mini-0853:phyxminiproblems-problemphyxmini0853-threeelectronescapesetup}
  Cartesian geometry for the aligned figure and the perturbed release state. The origin is the unperturbed centre-electron position. The actual release distances are metric distances, so the small displacement is not erased from the Coulomb-energy calculation
\end{definition}
\begin{proof}
  This is the declared model definition of Has Depicted And Release Geometry.
\end{proof}
\begin{definition}[Has Small Physical Perturbation]
  \label{def:physics:phyx-mini-0853:phyxminiproblems-problemphyxmini0853-hassmallphysicalperturbation}
  \lean{PhyXMiniProblems.ProblemPhyXMini0853.HasSmallPhysicalPerturbation}
  \uses{def:physics:phyx-mini-0853:phyxminiproblems-problemphyxmini0853-lengthinmeters, def:physics:phyx-mini-0853:phyxminiproblems-problemphyxmini0853-threeelectronescapesetup}
  The phrase "small distance" is quantified as at most one percent of the one-millimetre spacing. This is a conservative explicit idealization: it keeps the displacement nonzero while being far too small to change which widely separated answer choice is closest
\end{definition}
\begin{proof}
  This is the declared model definition of Has Small Physical Perturbation.
\end{proof}
\begin{definition}[Uses Electron And Vacuum Reference Data]
  \label{def:physics:phyx-mini-0853:phyxminiproblems-problemphyxmini0853-useselectronandvacuumreferencedata}
  \lean{PhyXMiniProblems.ProblemPhyXMini0853.UsesElectronAndVacuumReferenceData}
  \uses{def:physics:phyx-mini-0853:phyxminiproblems-problemphyxmini0853-massinkilograms, def:physics:phyx-mini-0853:phyxminiproblems-problemphyxmini0853-chargeincoulombs, def:physics:phyx-mini-0853:phyxminiproblems-problemphyxmini0853-coulombconstantinjoulemeterspercoulombsquared, def:physics:phyx-mini-0853:phyxminiproblems-problemphyxmini0853-elementarychargeincoulombs, def:physics:phyx-mini-0853:phyxminiproblems-problemphyxmini0853-threeelectronescapesetup}
  Standard electron and vacuum-electrostatics calibrations. These are reference data, not the requested far-away speed
\end{definition}
\begin{proof}
  This is the declared model definition of Uses Electron And Vacuum Reference Data.
\end{proof}
\begin{definition}[Satisfies Coulomb Escape And Energy Laws]
  \label{def:physics:phyx-mini-0853:phyxminiproblems-problemphyxmini0853-satisfiescoulombescapeandenergylaws}
  \lean{PhyXMiniProblems.ProblemPhyXMini0853.SatisfiesCoulombEscapeAndEnergyLaws}
  \uses{def:physics:phyx-mini-0853:phyxminiproblems-problemphyxmini0853-lengthinmeters, def:physics:phyx-mini-0853:phyxminiproblems-problemphyxmini0853-massinkilograms, def:physics:phyx-mini-0853:phyxminiproblems-problemphyxmini0853-chargeincoulombs, def:physics:phyx-mini-0853:phyxminiproblems-problemphyxmini0853-speedinmeterspersecond, def:physics:phyx-mini-0853:phyxminiproblems-problemphyxmini0853-energyinjoules, def:physics:phyx-mini-0853:phyxminiproblems-problemphyxmini0853-coulombconstantinjoulemeterspercoulombsquared, def:physics:phyx-mini-0853:phyxminiproblems-problemphyxmini0853-centerouterpair, def:physics:phyx-mini-0853:phyxminiproblems-problemphyxmini0853-threeelectronescapesetup}
  The centre electron's potential energy is the sum of its two pairwise Coulomb energies. At asymptotically infinite separations that potential vanishes. Kinetic energy obeys K = m v² / 2, and total mechanical energy is conserved. These laws are generic in the setup quantities and contain neither the target squared-speed formula nor any answer-choice value
\end{definition}
\begin{proof}
  This is the declared model definition of Satisfies Coulomb Escape And Energy Laws.
\end{proof}
\begin{definition}[Answer Choice]
  \label{def:physics:phyx-mini-0853:phyxminiproblems-problemphyxmini0853-answerchoice}
  \lean{PhyXMiniProblems.ProblemPhyXMini0853.AnswerChoice}
  Labels of the four speed choices printed with the problem
\end{definition}
\begin{proof}
  This is the declared model definition of Answer Choice.
\end{proof}
\begin{definition}[displayed Speed In Meters Per Second]
  \label{def:physics:phyx-mini-0853:phyxminiproblems-problemphyxmini0853-displayedspeedinmeterspersecond}
  \lean{PhyXMiniProblems.ProblemPhyXMini0853.displayedSpeedInMetersPerSecond}
  \uses{def:physics:phyx-mini-0853:phyxminiproblems-problemphyxmini0853-answerchoice}
  Speed printed beside each answer choice, in metres per second
\end{definition}
\begin{proof}
  This is the declared model definition of displayed Speed In Meters Per Second.
\end{proof}
\begin{definition}[Is Unique Closest Displayed Speed]
  \label{def:physics:phyx-mini-0853:phyxminiproblems-problemphyxmini0853-isuniqueclosestdisplayedspeed}
  \lean{PhyXMiniProblems.ProblemPhyXMini0853.IsUniqueClosestDisplayedSpeed}
  \uses{def:physics:phyx-mini-0853:phyxminiproblems-problemphyxmini0853-speedinmeterspersecond, def:physics:phyx-mini-0853:phyxminiproblems-problemphyxmini0853-threeelectronescapesetup, def:physics:phyx-mini-0853:phyxminiproblems-problemphyxmini0853-answerchoice, def:physics:phyx-mini-0853:phyxminiproblems-problemphyxmini0853-displayedspeedinmeterspersecond}
  A displayed choice is uniquely closest to the calculated far-away speed. This generic comparison predicate does not privilege any particular label
\end{definition}
\begin{proof}
  This is the declared model definition of Is Unique Closest Displayed Speed.
\end{proof}
% --- Archon declaration topology end ---
% --- Archon physics formalization source end ---
